\chapter{Matrix Basics: Notation, Addition, Transpose}

A matrix is a rectangular array of numbers organised into rows and columns.
In economics matrices store coefficients of linear equations, technology
coefficients in input-output models, covariances of asset returns, and design
matrices in regression. Before any deeper structure (determinants, inverses,
eigenvalues) can be discussed, one needs the basic vocabulary: dimensions,
entries, addition, scalar multiplication, and the transpose. This chapter
drills those mechanics.

\section{Core Rules}

\begin{rulebox}{Notation}
An $m \times n$ matrix $A \in \R^{m \times n}$ has $m$ rows and $n$ columns.
Its $(i, j)$-entry is written $a_{ij}$, so
\[
A \;=\; \bigl(a_{ij}\bigr)_{1 \le i \le m,\, 1 \le j \le n}.
\]
A row vector lies in $\R^{1 \times n}$, a column vector in $\R^{n \times 1}$
(or simply $\R^{n}$).
\end{rulebox}

\begin{rulebox}{Equality}
Two matrices are equal iff they have the same shape \emph{and} all entries
agree: $A = B$ iff $A, B \in \R^{m \times n}$ and $a_{ij} = b_{ij}$ for all
$i, j$.
\end{rulebox}

\begin{rulebox}{Addition and scalar multiplication}
For $A, B \in \R^{m \times n}$ and $\lambda \in \R$,
\[
(A + B)_{ij} = a_{ij} + b_{ij}, \qquad (\lambda A)_{ij} = \lambda \, a_{ij}.
\]
Addition requires matching dimensions. The set $\R^{m \times n}$ is a vector
space under these operations; in particular addition is commutative and
associative, and $\lambda(\mu A) = (\lambda \mu) A$.
\end{rulebox}

\begin{rulebox}{Zero, identity, diagonal}
The zero matrix $0$ has every entry $0$. The identity
$I_{n} \in \R^{n \times n}$ has $1$ on the diagonal and $0$ elsewhere. A
diagonal matrix has $a_{ij} = 0$ whenever $i \neq j$.
\end{rulebox}

\begin{rulebox}{Transpose}
The transpose $A^{T}$ of $A \in \R^{m \times n}$ is the $n \times m$ matrix
with $(A^{T})_{ij} = a_{ji}$. Rules:
\begin{itemize}[leftmargin=1.6em,itemsep=0.3ex]
  \item $(A^{T})^{T} = A$
  \item $(A + B)^{T} = A^{T} + B^{T}$
  \item $(\lambda A)^{T} = \lambda A^{T}$
\end{itemize}
\end{rulebox}

\begin{rulebox}{Symmetric and skew-symmetric}
A square matrix $A$ is \emph{symmetric} if $A^{T} = A$, equivalently
$a_{ij} = a_{ji}$. It is \emph{skew-symmetric} if $A^{T} = -A$ (in particular
all diagonal entries are zero). Covariance matrices, Hessians of $C^{2}$
functions, and the matrix $X^{T} X$ in regression are all symmetric.
\end{rulebox}

\section{Worked Example}

\begin{worked}{With $A=\begin{pmatrix}1&2&3\\4&5&6\end{pmatrix}$ and $B=\begin{pmatrix}1&0\\2&-1\\0&3\end{pmatrix}$, compute $2A-3B^{\!\top}$; then split $M=\begin{pmatrix}2&5\\1&4\end{pmatrix}$ into a symmetric and a skew-symmetric part}
\wstep{Step 1. Check the shapes before writing anything down.}
$A$ is $2\times3$ and $B$ is $3\times2$, so $A-B$ is \textbf{undefined}:
addition requires identical shapes. But $B^{\!\top}$ is $2\times3$, which
matches $A$. That is why the expression was written with a transpose, and
noticing this is the first step of the exercise.
\wtrap{Addition needs \emph{equal} dimensions; multiplication needs
\emph{matching} inner dimensions. These are different requirements, and the
first thing to do with any matrix expression is to verify that every
operation in it is legal.}

\wstep{Step 2. Transpose: rows become columns.}
\[
  B=\begin{pmatrix}1&0\\2&-1\\0&3\end{pmatrix}
  \quad\Longrightarrow\quad
  B^{\!\top}=\begin{pmatrix}1&2&0\\0&-1&3\end{pmatrix} .
\]
Entry $(i,j)$ of $B^{\!\top}$ is entry $(j,i)$ of $B$: the first
\emph{column} of $B$, namely $(1,2,0)$, becomes the first \emph{row}.

\wstep{Step 3. Scale each matrix, then subtract entry by entry.}
\[
  2A=\begin{pmatrix}2&4&6\\8&10&12\end{pmatrix},
  \qquad
  3B^{\!\top}=\begin{pmatrix}3&6&0\\0&-3&9\end{pmatrix},
\]
\[
  2A-3B^{\!\top}
  =\begin{pmatrix}2-3&4-6&6-0\\8-0&10-(-3)&12-9\end{pmatrix}
  =\boxed{\begin{pmatrix}-1&-2&6\\8&13&3\end{pmatrix}}
\]
\wtrap{$10-(-3)=13$. Most sign errors in this chapter come from subtracting
a negative entry. Scalar multiplication touches \emph{every} entry; a
frequent slip is to multiply only the first row.}

\wstep{Step 4. The rules that hold, and the one that does not.}
$(A+B)^{\!\top}=A^{\!\top}+B^{\!\top}$ and $(cA)^{\!\top}=cA^{\!\top}$:
transposition passes straight through sums and scalars. But for products the
order \textbf{reverses}:
\[
  (AB)^{\!\top}=B^{\!\top}A^{\!\top}\neq A^{\!\top}B^{\!\top}
\]
in general --- indeed $A^{\!\top}B^{\!\top}$ is usually not even defined.
Check the shapes: $A$ is $2\times3$, $B$ is $3\times2$, so $AB$ is
$2\times2$ and $(AB)^{\!\top}$ is $2\times2$; $B^{\!\top}A^{\!\top}$ is
$(2\times3)(3\times2)=2\times2$ \checkmark, while $A^{\!\top}B^{\!\top}$
would be $(3\times2)(2\times3)=3\times3$. The shapes alone prove which
formula is right.

\wstep{Step 5. The symmetric--skew split, which every square matrix admits.}
\[
  M=\underbrace{\tfrac12(M+M^{\!\top})}_{\text{symmetric}}
   +\underbrace{\tfrac12(M-M^{\!\top})}_{\text{skew-symmetric}} .
\]
With $M=\begin{pmatrix}2&5\\1&4\end{pmatrix}$ and
$M^{\!\top}=\begin{pmatrix}2&1\\5&4\end{pmatrix}$:
\[
  \tfrac12(M+M^{\!\top})=\tfrac12\begin{pmatrix}4&6\\6&8\end{pmatrix}
  =\begin{pmatrix}2&3\\3&4\end{pmatrix},
  \qquad
  \tfrac12(M-M^{\!\top})=\tfrac12\begin{pmatrix}0&4\\-4&0\end{pmatrix}
  =\begin{pmatrix}0&2\\-2&0\end{pmatrix} .
\]
Their sum is $\begin{pmatrix}2&5\\1&4\end{pmatrix}=M$. \checkmark

\wstep{Step 6. Why the split matters later.}
The first part is symmetric ($S^{\!\top}=S$), the second skew
($K^{\!\top}=-K$, so its diagonal is forced to be zero). A quadratic form
$\mathbf x^{\!\top}M\mathbf x$ sees only the symmetric part, because
$\mathbf x^{\!\top}K\mathbf x=0$ for every skew $K$. That is why every
Hessian in \volref{derivatives} and every definiteness test in \chref{matrices}{9} is stated for a
\emph{symmetric} matrix: nothing is lost, and the eigenvalues are then
guaranteed real.
\wtrap{Symmetry is about the entries, not about being square. A square
matrix need not be symmetric, and only a square matrix can be. Checking
$M=M^{\!\top}$ means checking $m_{ij}=m_{ji}$ for every off-diagonal pair
--- one glance at $5$ against $1$ settles it here.}
\end{worked}

\section{Practice Problems}

Use the following matrices throughout:
\[
A = \begin{pmatrix} 1 & 2 \\ 3 & 4 \end{pmatrix}, \quad
B = \begin{pmatrix} 0 & -1 \\ 2 & 5 \end{pmatrix}, \quad
C = \begin{pmatrix} 1 & 0 & 2 \\ -1 & 3 & 1 \end{pmatrix}, \quad
D = \begin{pmatrix} 2 & 1 & 0 \\ 0 & 4 & -1 \end{pmatrix}.
\]

\begin{problems}
  \item State the dimensions of $A$, $B$, $C$, $D$.
  \item Compute $A + B$.
  \item Compute $A - B$.
  \item Compute $3 A$.
  \item Compute $2 A - 3 B$.
  \item Compute $C + D$.
  \item Compute $C - 2 D$.
  \item Compute $A^{T}$.
  \item Compute $B^{T}$.
  \item Compute $C^{T}$.
  \item Compute $(A + B)^{T}$ and verify it equals $A^{T} + B^{T}$.
  \item Is the matrix $\bigl(\begin{smallmatrix} 1 & 2 \\ 2 & 5 \end{smallmatrix}\bigr)$ symmetric? Skew-symmetric? Neither?
  \item Is the matrix $\bigl(\begin{smallmatrix} 0 & 3 \\ -3 & 0 \end{smallmatrix}\bigr)$ symmetric? Skew-symmetric? Neither?
  \item Write the $3 \times 3$ identity matrix $I_{3}$.
  \item Find all $x$ such that $\bigl(\begin{smallmatrix} x & 2 \\ 2 & x-1 \end{smallmatrix}\bigr)$ is symmetric.
  \item Find all $x, y$ such that $\bigl(\begin{smallmatrix} 0 & x \\ y & 0 \end{smallmatrix}\bigr)$ is skew-symmetric.
  \item Compute $A + A^{T}$. Is it symmetric?
  \item Compute $A - A^{T}$. Is it skew-symmetric?
  \item Solve for the matrix $X$: $X + A = B$.
  \item Solve for the matrix $X$: $2 X - 3 A = B$.
  \item Can you compute $A + C$? Justify briefly.
  \item Write the $2 \times 2$ zero matrix and verify $A + 0 = A$.
  \item Show by computation that $(2 A)^{T} = 2 A^{T}$.
  \item Compute the trace $\mathrm{tr}(A) = a_{11} + a_{22}$ for $A$ above, and for $B$.
  \item A diagonal matrix $\mathrm{diag}(2, -1, 3)$ has what dimensions and what $(2,2)$-entry?
\end{problems}

\section{Solutions}

\begin{solutions}
  \item $A, B$ are $2 \times 2$. $C, D$ are $2 \times 3$.

  \item $A + B = \bigl(\begin{smallmatrix} 1 & 1 \\ 5 & 9 \end{smallmatrix}\bigr)$.

  \item $A - B = \bigl(\begin{smallmatrix} 1 & 3 \\ 1 & -1 \end{smallmatrix}\bigr)$.

  \item $3 A = \bigl(\begin{smallmatrix} 3 & 6 \\ 9 & 12 \end{smallmatrix}\bigr)$.

  \item $2 A - 3 B = \bigl(\begin{smallmatrix} 2 & 4 \\ 6 & 8 \end{smallmatrix}\bigr) - \bigl(\begin{smallmatrix} 0 & -3 \\ 6 & 15 \end{smallmatrix}\bigr) = \bigl(\begin{smallmatrix} 2 & 7 \\ 0 & -7 \end{smallmatrix}\bigr)$.

  \item $C + D = \bigl(\begin{smallmatrix} 3 & 1 & 2 \\ -1 & 7 & 0 \end{smallmatrix}\bigr)$.

  \item $C - 2 D = \bigl(\begin{smallmatrix} 1 & 0 & 2 \\ -1 & 3 & 1 \end{smallmatrix}\bigr) - \bigl(\begin{smallmatrix} 4 & 2 & 0 \\ 0 & 8 & -2 \end{smallmatrix}\bigr) = \bigl(\begin{smallmatrix} -3 & -2 & 2 \\ -1 & -5 & 3 \end{smallmatrix}\bigr)$.

  \item $A^{T} = \bigl(\begin{smallmatrix} 1 & 3 \\ 2 & 4 \end{smallmatrix}\bigr)$.

  \item $B^{T} = \bigl(\begin{smallmatrix} 0 & 2 \\ -1 & 5 \end{smallmatrix}\bigr)$.

  \item $C^{T} = \bigl(\begin{smallmatrix} 1 & -1 \\ 0 & 3 \\ 2 & 1 \end{smallmatrix}\bigr) \in \R^{3 \times 2}$.

  \item $(A + B)^{T} = \bigl(\begin{smallmatrix} 1 & 1 \\ 5 & 9 \end{smallmatrix}\bigr)^{T} = \bigl(\begin{smallmatrix} 1 & 5 \\ 1 & 9 \end{smallmatrix}\bigr)$, and $A^{T} + B^{T} = \bigl(\begin{smallmatrix} 1 & 3 \\ 2 & 4 \end{smallmatrix}\bigr) + \bigl(\begin{smallmatrix} 0 & 2 \\ -1 & 5 \end{smallmatrix}\bigr) = \bigl(\begin{smallmatrix} 1 & 5 \\ 1 & 9 \end{smallmatrix}\bigr)$. Equal.

  \item Symmetric, since $a_{12} = a_{21} = 2$.

  \item Skew-symmetric: $a_{12} = 3$, $a_{21} = -3 = -a_{12}$, and diagonal entries are $0$.

  \item $I_{3} = \bigl(\begin{smallmatrix} 1 & 0 & 0 \\ 0 & 1 & 0 \\ 0 & 0 & 1 \end{smallmatrix}\bigr)$.

  \item Off-diagonal entries are both $2$ regardless of $x$, so the matrix is symmetric for \emph{every} $x \in \R$.

  \item Skew-symmetry needs $a_{21} = -a_{12}$, i.e.\ $y = -x$. So $\{(x, -x) : x \in \R\}$.

  \item $A + A^{T} = \bigl(\begin{smallmatrix} 1 & 2 \\ 3 & 4 \end{smallmatrix}\bigr) + \bigl(\begin{smallmatrix} 1 & 3 \\ 2 & 4 \end{smallmatrix}\bigr) = \bigl(\begin{smallmatrix} 2 & 5 \\ 5 & 8 \end{smallmatrix}\bigr)$, symmetric (general fact: $A + A^{T}$ is always symmetric).

  \item $A - A^{T} = \bigl(\begin{smallmatrix} 0 & -1 \\ 1 & 0 \end{smallmatrix}\bigr)$, skew-symmetric.

  \item $X = B - A = \bigl(\begin{smallmatrix} -1 & -3 \\ -1 & 1 \end{smallmatrix}\bigr)$.

  \item $2 X = B + 3 A = \bigl(\begin{smallmatrix} 0 & -1 \\ 2 & 5 \end{smallmatrix}\bigr) + \bigl(\begin{smallmatrix} 3 & 6 \\ 9 & 12 \end{smallmatrix}\bigr) = \bigl(\begin{smallmatrix} 3 & 5 \\ 11 & 17 \end{smallmatrix}\bigr)$, so $X = \bigl(\begin{smallmatrix} 3/2 & 5/2 \\ 11/2 & 17/2 \end{smallmatrix}\bigr)$.

  \item No: $A$ is $2 \times 2$ and $C$ is $2 \times 3$, dimensions disagree.

  \item $0 = \bigl(\begin{smallmatrix} 0 & 0 \\ 0 & 0 \end{smallmatrix}\bigr)$, and entry-wise $A + 0 = A$.

  \item $2 A = \bigl(\begin{smallmatrix} 2 & 4 \\ 6 & 8 \end{smallmatrix}\bigr)$, so $(2 A)^{T} = \bigl(\begin{smallmatrix} 2 & 6 \\ 4 & 8 \end{smallmatrix}\bigr)$. Also $2 A^{T} = 2 \bigl(\begin{smallmatrix} 1 & 3 \\ 2 & 4 \end{smallmatrix}\bigr) = \bigl(\begin{smallmatrix} 2 & 6 \\ 4 & 8 \end{smallmatrix}\bigr)$. Equal.

  \item $\mathrm{tr}(A) = 1 + 4 = 5$, $\mathrm{tr}(B) = 0 + 5 = 5$.

  \item $\mathrm{diag}(2, -1, 3) \in \R^{3 \times 3}$ and its $(2,2)$-entry is $-1$.
\end{solutions}

\section{Additional Problems}

Use the following matrices throughout:
\[
E = \begin{pmatrix} 2 & -1 \\ 0 & 3 \end{pmatrix}, \quad
F = \begin{pmatrix} 1 & 4 \\ -2 & 2 \end{pmatrix}, \quad
G = \begin{pmatrix} 3 & 0 & -1 \\ 2 & 1 & 4 \end{pmatrix}, \quad
H = \begin{pmatrix} -1 & 2 & 5 \\ 0 & 3 & 1 \end{pmatrix}.
\]

\begin{problems}
  \item Compute $E + F$.
  \item Compute $E - F$.
  \item Compute $4 E$.
  \item Compute $3 E + 2 F$.
  \item Compute $G - H$.
  \item Compute $2 G + H$.
  \item Compute $E^{T}$.
  \item Compute $G^{T}$.
  \item Verify $(E + F)^{T} = E^{T} + F^{T}$.
  \item Is $\bigl(\begin{smallmatrix} 3 & -2 \\ -2 & 7 \end{smallmatrix}\bigr)$ symmetric, skew-symmetric, or neither?
  \item Is $\bigl(\begin{smallmatrix} 0 & 5 & -2 \\ -5 & 0 & 1 \\ 2 & -1 & 0 \end{smallmatrix}\bigr)$ symmetric, skew-symmetric, or neither?
  \item Find all $a$ such that $\bigl(\begin{smallmatrix} 2 & a+1 \\ 3a-5 & 7 \end{smallmatrix}\bigr)$ is symmetric.
  \item Solve for $X$: $2 X - E = 3 F$.
  \item Solve for $X$: $X^{T} = E + F$.
  \item Compute $\mathrm{tr}(E)$ and $\mathrm{tr}(F)$, and verify $\mathrm{tr}(E + F) = \mathrm{tr}(E) + \mathrm{tr}(F)$.
  \item Decompose $E$ as a sum of a symmetric and a skew-symmetric matrix: $E = \tfrac{1}{2}(E + E^{T}) + \tfrac{1}{2}(E - E^{T})$.
  \item For $E$, compute the matrix $E + E^{T}$ and verify it is symmetric.
  \item Can you compute $G + E$? Justify.
  \item Find all $b, c$ such that $\bigl(\begin{smallmatrix} 0 & b & 1 \\ c & 0 & -3 \\ -1 & 3 & 0 \end{smallmatrix}\bigr)$ is skew-symmetric.
  \item Compute the $(2, 1)$-entry of $E^{T}$ and confirm it equals the $(1, 2)$-entry of $E$.
  \item Write the $2 \times 3$ zero matrix and check $G + 0 = G$.
  \item For $D = \mathrm{diag}(5, -2)$, compute $D + D^{T}$.
  \item Show by computation that $(-3 F)^{T} = -3 F^{T}$.
  \item For the matrix $\bigl(\begin{smallmatrix} 4 & 2 \\ 2 & 4 \end{smallmatrix}\bigr)$, find its trace and check whether it is symmetric.
  \item A skew-symmetric $3 \times 3$ matrix has how many free parameters?
\end{problems}

\section{Additional Solutions}

\begin{solutions}
  \item $E + F = \bigl(\begin{smallmatrix} 3 & 3 \\ -2 & 5 \end{smallmatrix}\bigr)$.

  \item $E - F = \bigl(\begin{smallmatrix} 1 & -5 \\ 2 & 1 \end{smallmatrix}\bigr)$.

  \item $4 E = \bigl(\begin{smallmatrix} 8 & -4 \\ 0 & 12 \end{smallmatrix}\bigr)$.

  \item $3 E + 2 F = \bigl(\begin{smallmatrix} 6 & -3 \\ 0 & 9 \end{smallmatrix}\bigr) + \bigl(\begin{smallmatrix} 2 & 8 \\ -4 & 4 \end{smallmatrix}\bigr) = \bigl(\begin{smallmatrix} 8 & 5 \\ -4 & 13 \end{smallmatrix}\bigr)$.

  \item $G - H = \bigl(\begin{smallmatrix} 4 & -2 & -6 \\ 2 & -2 & 3 \end{smallmatrix}\bigr)$.

  \item $2 G + H = \bigl(\begin{smallmatrix} 6 & 0 & -2 \\ 4 & 2 & 8 \end{smallmatrix}\bigr) + \bigl(\begin{smallmatrix} -1 & 2 & 5 \\ 0 & 3 & 1 \end{smallmatrix}\bigr) = \bigl(\begin{smallmatrix} 5 & 2 & 3 \\ 4 & 5 & 9 \end{smallmatrix}\bigr)$.

  \item $E^{T} = \bigl(\begin{smallmatrix} 2 & 0 \\ -1 & 3 \end{smallmatrix}\bigr)$.

  \item $G^{T} = \bigl(\begin{smallmatrix} 3 & 2 \\ 0 & 1 \\ -1 & 4 \end{smallmatrix}\bigr) \in \R^{3 \times 2}$.

  \item $(E + F)^{T} = \bigl(\begin{smallmatrix} 3 & 3 \\ -2 & 5 \end{smallmatrix}\bigr)^{T} = \bigl(\begin{smallmatrix} 3 & -2 \\ 3 & 5 \end{smallmatrix}\bigr)$, and $E^{T} + F^{T} = \bigl(\begin{smallmatrix} 2 & 0 \\ -1 & 3 \end{smallmatrix}\bigr) + \bigl(\begin{smallmatrix} 1 & -2 \\ 4 & 2 \end{smallmatrix}\bigr) = \bigl(\begin{smallmatrix} 3 & -2 \\ 3 & 5 \end{smallmatrix}\bigr)$. Equal.

  \item Symmetric: $a_{12} = a_{21} = -2$.

  \item Skew-symmetric: all $a_{ij} = -a_{ji}$ and diagonal entries are $0$.

  \item Need $a + 1 = 3 a - 5 \Rightarrow 2 a = 6 \Rightarrow a = 3$.

  \item $2 X = E + 3 F = \bigl(\begin{smallmatrix} 2 & -1 \\ 0 & 3 \end{smallmatrix}\bigr) + \bigl(\begin{smallmatrix} 3 & 12 \\ -6 & 6 \end{smallmatrix}\bigr) = \bigl(\begin{smallmatrix} 5 & 11 \\ -6 & 9 \end{smallmatrix}\bigr)$, so $X = \bigl(\begin{smallmatrix} 5/2 & 11/2 \\ -3 & 9/2 \end{smallmatrix}\bigr)$.

  \item $X = (E + F)^{T} = \bigl(\begin{smallmatrix} 3 & -2 \\ 3 & 5 \end{smallmatrix}\bigr)$.

  \item $\mathrm{tr}(E) = 2 + 3 = 5$, $\mathrm{tr}(F) = 1 + 2 = 3$. $\mathrm{tr}(E + F) = 3 + 5 = 8 = 5 + 3$. Check.

  \item Symmetric part: $\tfrac{1}{2}(E + E^{T}) = \tfrac{1}{2}\bigl(\begin{smallmatrix} 4 & -1 \\ -1 & 6 \end{smallmatrix}\bigr) = \bigl(\begin{smallmatrix} 2 & -1/2 \\ -1/2 & 3 \end{smallmatrix}\bigr)$. Skew part: $\tfrac{1}{2}(E - E^{T}) = \tfrac{1}{2}\bigl(\begin{smallmatrix} 0 & -1 \\ 1 & 0 \end{smallmatrix}\bigr) = \bigl(\begin{smallmatrix} 0 & -1/2 \\ 1/2 & 0 \end{smallmatrix}\bigr)$. Sum returns $E$.

  \item $E + E^{T} = \bigl(\begin{smallmatrix} 4 & -1 \\ -1 & 6 \end{smallmatrix}\bigr)$. The off-diagonal entries are both $-1$: symmetric.

  \item No: $G$ is $2 \times 3$ and $E$ is $2 \times 2$, dimensions disagree.

  \item Skew-symmetry requires $a_{12} = -a_{21}$ and $a_{13} = -a_{31}$, $a_{23} = -a_{32}$. Here $a_{13} = 1, a_{31} = -1$: OK. $a_{23} = -3, a_{32} = 3$: OK. So need $b = -c$, i.e.\ $(b, c) = (t, -t)$ for any $t \in \R$.

  \item $(E^{T})_{21} = E_{12} = -1$. Confirmed.

  \item $0 = \bigl(\begin{smallmatrix} 0 & 0 & 0 \\ 0 & 0 & 0 \end{smallmatrix}\bigr)$ and $G + 0 = G$ entry-wise.

  \item $D + D^{T} = 2 D = \bigl(\begin{smallmatrix} 10 & 0 \\ 0 & -4 \end{smallmatrix}\bigr)$ (since $D$ is diagonal, hence symmetric).

  \item $-3 F = \bigl(\begin{smallmatrix} -3 & -12 \\ 6 & -6 \end{smallmatrix}\bigr)$, so $(-3 F)^{T} = \bigl(\begin{smallmatrix} -3 & 6 \\ -12 & -6 \end{smallmatrix}\bigr)$. Also $-3 F^{T} = -3 \bigl(\begin{smallmatrix} 1 & -2 \\ 4 & 2 \end{smallmatrix}\bigr) = \bigl(\begin{smallmatrix} -3 & 6 \\ -12 & -6 \end{smallmatrix}\bigr)$. Equal.

  \item $\mathrm{tr} = 4 + 4 = 8$. The matrix is symmetric since the off-diagonal entries are both $2$.

  \item A $3 \times 3$ skew-symmetric matrix has zero diagonal and the entries below the diagonal are minus those above; only the entries strictly above the diagonal are free. That gives $3$ free parameters ($a_{12}, a_{13}, a_{23}$).
\end{solutions}
